\begin{table}[H]
\centering
\caption{Hyperparameter validation of the R2 degree null (TransE and DistMult; sampled 50 type-matched negatives, the protocol of Table~\ref{tab:audit}). \textbf{Panel A} scores roughly 4,000 held-\emph{in} training edges of the target relation through the same evaluator that produces the test numbers, so the only difference between the train-fit and test columns is whether the model saw the edge during training. Train fit upper-bounds what the optimizer achieved on the graph it was given: a recipe unable to fit the rewired graph would show train fit collapsing with test fit, and instead both models nearly preserve train fit while losing a quarter to a third of held-out performance, which is a generalization failure rather than an optimization failure. \textbf{Panel B} runs one learning-rate $\times$ negatives-per-positive grid on R2 and on R0, the R0 arm serving as the matched control. Panel B verdict: the best R0 configuration (d64 e100 lr=0.001 neg=64, MRR 0.8167) and the best R2 configuration (d64 e100 lr=0.001 neg=64, MRR 0.6121) are the SAME configuration, so R2 is not handicapped by having inherited tuning chosen on R0 -- which contradicts the objection's premise directly. Best-vs-best gap -25.0\% (R2 retains 0.750); R2 trails R0 in 6 of 6 paired cells, with the narrowest gap anywhere in the grid at -23.6\%. The grid was capable of moving performance (R0 spans 0.2327 MRR, R2 0.2106), so its failure to close the gap is informative rather than a dead grid. Panel B runs 100 epochs rather than the headline 300 because it is asked for configuration ORDERING, not headline numbers, and its cells should not be quoted as results; the cost of that is measurable rather than assumed, since the lr 1e-3 / neg 16 cell is exactly the headline recipe and differs from the frozen 300-epoch seed-42 run only in epoch count -- it lands +0.2\% from it on R0 and +0.2\% on R2, so TransE has effectively converged by 100 epochs on both graphs and the R2 arm is not penalised by the short budget. Thresholds were fixed before the runs: train-fit retention $\geq$ 0.85, fit-vs-test separation $\geq$ 0.15, and a surviving best-vs-best gap $\geq$ 15\%. The embedding-dimension axis was not swept, so this grid varies optimization hyperparameters and not model capacity; a capacity shortfall would appear as depressed R2 train fit, which Panel A measures directly.}
\label{tab:r2_sweep}
\begin{tabular}{lrrrrrrrr}
\toprule
Model / config & R0 MRR & R2 MRR & $\Delta$MRR R0$\to$R2 & R0 train-fit MRR & R2 train-fit MRR & Train-fit retention & Test retention & Separation \\
\midrule
\multicolumn{9}{l}{\textit{Panel A. Train-fit diagnostic (300 epochs, seed 42, headline recipe)}} \\
DistMult & 0.6665 & 0.4208 & -36.9\% & 0.8842 & 0.7720 & 0.873 & 0.631 & 0.242 \\
TransE & 0.7975 & 0.6090 & -23.6\% & 0.9120 & 0.8565 & 0.939 & 0.764 & 0.175 \\
\multicolumn{9}{l}{\textit{Panel B. Matched hyperparameter grid, same grid on both regimes (100 epochs, seed 42)}} \\
d64 e100 lr=0.001 neg=16 & 0.7988 & 0.6101 & -23.6\% & -- & -- & -- & 0.764 & -- \\
d64 e100 lr=0.001 neg=64 & 0.8167 & 0.6121 & -25.0\% & -- & -- & -- & 0.750 & -- \\
d64 e100 lr=0.003 neg=16 & 0.7517 & 0.5470 & -27.2\% & -- & -- & -- & 0.728 & -- \\
d64 e100 lr=0.003 neg=64 & 0.7798 & 0.5647 & -27.6\% & -- & -- & -- & 0.724 & -- \\
d64 e100 lr=0.01 neg=16 & 0.5840 & 0.4015 & -31.2\% & -- & -- & -- & 0.688 & -- \\
d64 e100 lr=0.01 neg=64 & 0.6428 & 0.4130 & -35.8\% & -- & -- & -- & 0.642 & -- \\
\bottomrule
\end{tabular}
\end{table}
